% ============================================================================
%  08_herramientas.tex — IEEE 730 §8: Herramientas (Sección A.c)
% ============================================================================

\chapter{Herramientas, Técnicas y Metodologías}
\label{ch:herramientas}

La selección de herramientas SQA respondió a tres criterios: \textbf{compatibilidad} con el stack
tecnológico del proyecto (Python/FastAPI, Next.js), \textbf{disponibilidad} en el entorno de pruebas
(sin costo de licencia en el nivel académico) y \textbf{cobertura} de las diferentes capas de calidad
(estática, unitaria, integración, sistema, CI/CD).

% ---------------------------------------------------------------------------
\section{Herramientas de Gestión de Calidad}
\label{sec:herramientas-gestion}

\subsection{Jira Software (Gestión de Defectos y Tareas)}
\label{subsec:jira}

\textbf{Rol en el proyecto:} Repositorio oficial de defectos y tablero de seguimiento del sprint SQA.

\textbf{Justificación técnica:} Jira es el estándar de facto en gestión ágil de proyectos de software y
ofrece integración nativa con GitHub, lo que permite vincular commits y pull requests directamente a tickets
de defecto. En el contexto de este proyecto, se utilizó para:

\begin{itemize}
    \item Registrar los 9 defectos encontrados durante el sprint (BUG-001 a BUG-009).
    \item Clasificar defectos por severidad (Critical, Major, Minor) y prioridad.
    \item Hacer seguimiento del ciclo de vida de cada defecto (Abierto → En Progreso → Resuelto → Cerrado).
    \item Generar el informe de densidad de defectos del Apéndice C.
\end{itemize}

\textbf{Alternativas evaluadas:} GitHub Issues (descartado por falta de campos de severidad/prioridad
nativos), Trello (descartado por falta de flujo de estados).

\subsection{SonarCloud (Análisis Estático de Código)}
\label{subsec:sonarcloud}

\textbf{Rol en el proyecto:} Análisis continuo de calidad de código para el repositorio completo.

\textbf{Justificación técnica:} SonarCloud es la solución en la nube de SonarQube, compatible de forma nativa
con Python (backend FastAPI) y TypeScript/JavaScript (frontend Next.js). Proporciona métricas objetivas de:

\begin{itemize}
    \item \textbf{Deuda técnica:} tiempo estimado para corregir todos los \textit{code smells}.
    \item \textbf{Cobertura de código:} porcentaje de líneas cubiertas por pruebas automatizadas.
    \item \textbf{Duplicación:} porcentaje de código duplicado en el repositorio.
    \item \textbf{Vulnerabilidades de seguridad:} detección de patrones de OWASP Top 10.
    \item \textbf{Reliability:} bugs potenciales detectados por análisis de flujo de datos.
\end{itemize}

\textbf{Alternativas evaluadas:} Pylint + ESLint (herramientas locales, descartadas porque no
proporcionan un portal centralizado ni integración con CI/CD de forma inmediata).

% ---------------------------------------------------------------------------
\section{Herramientas de Automatización de Pruebas}
\label{sec:herramientas-automatizacion}

\subsection{pytest (Pruebas Unitarias e Integración del Backend)}
\label{subsec:pytest}

\textbf{Rol en el proyecto:} Framework de pruebas para el backend Python/FastAPI.

\textbf{Justificación técnica:} pytest es la herramienta estándar del ecosistema Python. Se eligió sobre
\texttt{unittest} (biblioteca estándar) por:

\begin{itemize}
    \item Sintaxis más concisa (sin clases ni métodos \texttt{setUp}).
    \item Soporte nativo de \textit{fixtures} para inyección de dependencias en tests.
    \item Compatible con \texttt{httpx.AsyncClient} para pruebas de endpoints FastAPI de forma asíncrona.
    \item Integración directa con SonarCloud para reporte de cobertura vía \texttt{pytest-cov}.
\end{itemize}

\textbf{Módulos cubiertos:} Auth (\texttt{test\_auth.py}), Transcripciones (\texttt{test\_transcriptions.py}),
Health (\texttt{test\_health.py}), VAD Config (\texttt{test\_vad\_config.py}), Seguridad de Rutas
(\texttt{test\_path\_security.py}), Settings (\texttt{test\_settings.py}).

\subsection{Cypress (Pruebas E2E del Frontend)}
\label{subsec:cypress}

\textbf{Rol en el proyecto:} Automatización de pruebas de extremo a extremo sobre el frontend Next.js.

\textbf{Justificación técnica:} Cypress es la herramienta E2E más adoptada en el ecosistema JavaScript/TypeScript.
Se eligió sobre Playwright por su \textit{Test Runner} visual y la facilidad de depuración interactiva. Permite
automatizar el flujo completo de usuario: login → grabación → transcripción → edición, ejecutando las pruebas
directamente en el navegador Chromium sin necesidad de drivers externos (Selenium WebDriver).

\subsection{Jest (Pruebas Unitarias del Frontend)}
\label{subsec:jest}

\textbf{Rol en el proyecto:} Pruebas unitarias de componentes React en el frontend Next.js.

\textbf{Justificación técnica:} Jest es el framework de pruebas por defecto de Next.js (incluido en su
configuración de boilerplate). Junto con \texttt{@testing-library/react}, permite verificar el renderizado
y el comportamiento de componentes como \texttt{RecordingPanel} y \texttt{LiveTranscription} sin necesidad
de levantar un servidor.

\subsection{Postman / curl (Pruebas Manuales de API)}
\label{subsec:postman}

\textbf{Rol en el proyecto:} Pruebas exploratorias y de contrato sobre los endpoints REST del backend.

\textbf{Justificación técnica:} Postman proporciona una interfaz gráfica para diseñar y ejecutar solicitudes
HTTP, documentar contratos de API y compartir colecciones entre los miembros del equipo SQA. Se utilizó
curl como alternativa en consola para pruebas rápidas y para el script de verificación del endpoint
\texttt{/health}. No se requirió instalación adicional en el entorno de pruebas.

% ---------------------------------------------------------------------------
\section{Herramientas de Integración Continua}
\label{sec:herramientas-ci}

\subsection{GitHub Actions (CI/CD)}
\label{subsec:github-actions}

\textbf{Rol en el proyecto:} Pipeline de integración continua que ejecuta automáticamente las pruebas
unitarias e integración del backend en cada \textit{pull request} hacia la rama \texttt{main}.

\textbf{Justificación técnica:} GitHub Actions es la solución de CI/CD nativa de GitHub, sin costo para
repositorios públicos. El workflow \texttt{.github/workflows/} configurado en el proyecto ejecuta:

\begin{enumerate}
    \item Instalación de dependencias Python (\texttt{pip install -r requirements.txt}).
    \item Ejecución de la suite pytest con reporte de cobertura.
    \item Publicación del reporte de cobertura a SonarCloud.
\end{enumerate}

\textbf{Alternativas evaluadas:} Jenkins (requiere servidor propio), CircleCI (requiere configuración de
créditos de cómputo).

% ---------------------------------------------------------------------------
\section{Resumen del Stack de Herramientas SQA}
\label{sec:resumen-herramientas}

\begin{longtable}{p{3cm}p{1.8cm}p{2.5cm}p{2.5cm}p{3cm}}
\caption{Resumen de herramientas SQA seleccionadas}
\label{tab:herramientas-resumen}\\
\toprule
\textbf{Herramienta} & \textbf{Versión} & \textbf{Categoría} & \textbf{Capa} & \textbf{Propósito}\\
\midrule
\endfirsthead
\toprule
\textbf{Herramienta} & \textbf{Versión} & \textbf{Categoría} & \textbf{Capa} & \textbf{Propósito}\\
\midrule
\endhead
\bottomrule
\endfoot

Jira Software & Cloud & Gestión & Proceso & Seguimiento de defectos y tareas del sprint\\
SonarCloud & Cloud & Análisis Estático & Código & Deuda técnica, vulnerabilidades, cobertura\\
GitHub Actions & Cloud & CI/CD & Proceso & Ejecución automática de pruebas en PRs\\
pytest & 8.x & Automatización & Backend & Pruebas unitarias e integración (Python)\\
Cypress & 14.x & Automatización & Frontend/E2E & Pruebas de flujo de usuario completo\\
Jest & 29.x & Automatización & Frontend & Pruebas unitarias de componentes React\\
Postman & Desktop & Manual & API & Pruebas exploratorias de endpoints REST\\
curl & CLI & Manual & API & Verificación rápida de endpoints REST\\

\bottomrule
\end{longtable}
